\chapter{Αποτελέσματα και Ανάλυση Δεδομένων}

\section{Αποτελέσματα Ακουστικής Άνεσης στο Μαθησιακό Περιβάλλον}
Η εφαρμογή του Ψηφιακού Διδύμου επέφερε άμεσα και μετρήσιμα αποτελέσματα ως προς την Ποιότητα Εσωτερικού Περιβάλλοντος (\textlatin{IEQ}). Η συλλογή των αρχικών τηλεμετρικών δεδομένων για τον θόρυβο πραγματοποιήθηκε κατά τη φάση του τοπικού \textlatin{MVP}, μέσω της σύζευξης \textlatin{ESP32}, \textlatin{MQTT}, \textlatin{openHAB} και \textlatin{InfluxDB}, η οποία επέτρεψε την αξιόπιστη αποτύπωση των υπερβάσεων του επιχειρησιακού ορίου των 55 \textlatin{dBA}. Στη συνέχεια, η νέα \textlatin{Enterprise Cloud SaaS} πλατφόρμα αξιοποιήθηκε για την ενοποιημένη ανάλυση, την πιο φιλική οπτικοποίηση και την ευανάγνωστη παρουσίαση των αποτελεσμάτων προς τεχνικούς και διοικητικούς χρήστες.

Η χρήση του κλειστού βρόχου ελέγχου για τη διατήρηση του θορύβου κάτω από το δυναμικό όριο των 55 \textlatin{dBA} αποδείχθηκε ιδιαίτερα αποτελεσματική. Σε αντίθεση με τη χρήση ηχητικών συναγερμών (\textlatin{buzzers}), η οποία θα διατάρασσε την εκπαιδευτική διαδικασία, η διακριτική οπτική ανατροφοδότηση (\textlatin{quiet-lamp}) λειτούργησε ομαλά ως μηχανισμός αυτορρύθμισης της συμπεριφοράς των μαθητών. Τα αρχικά δεδομένα από την υποδομή παρακολούθησης και η μεταγενέστερη ενοποιημένη απεικόνισή τους στο \textlatin{cloud} περιβάλλον, όπως αποτυπώνεται στο Σχήμα~\ref{fig:noise_results}, κατέδειξαν σημαντική μείωση στον αριθμό των ημερήσιων υπερβάσεων του ορίου, συμβάλλοντας στη διαμόρφωση περιβάλλοντος περισσότερο συμβατού με τις κατευθυντήριες οδηγίες ακουστικής άνεσης.

\begin{figure}[htbp]
\centering
\includegraphics[width=\textwidth]{fig_6_1}
\caption{Απεικόνιση επιπέδων θορύβου και υπερβάσεων (\textlatin{exceedances}) εντός της σχολικής αίθουσας}
\label{fig:noise_results}
\end{figure}

Αναφορικά με την ακρίβεια των αισθητήρων, οι κόμβοι \textlatin{ESP32/INMP441} επαληθεύτηκαν έναντι πιστοποιημένου βιομηχανικού ηχομέτρου αναφοράς, επιτυγχάνοντας μέση ακρίβεια ±2 \textlatin{dBA}. Αντίστοιχα, οι αισθητήρες \textlatin{DHT22} για θερμοκρασία και υγρασία επιβεβαιώθηκαν με ακρίβεια ±0.5°C και ±2\% \textlatin{RH} αντίστοιχα μετά από βαθμονόμηση. Η συνολική αξιοπιστία δεδομένων (\textlatin{data quality score}) του συστήματος ανήλθε σε 98.5\%, τιμή που υπολογίστηκε μέσω διασταύρωσης των μετρήσεων των αισθητήρων με ανεξάρτητα όργανα αναφοράς.

\section{Ανάλυση Ενεργειακής Απόδοσης και Σύγκριση Δεδομένων}
Για την αξιολόγηση της ενεργειακής επίδρασης της πλατφόρμας, πραγματοποιήθηκε σύγκριση της μέσης ημερήσιας κατανάλωσης (\textlatin{kWh}/ημέρα) μεταξύ της περιόδου πριν και μετά την εφαρμογή του Ψηφιακού Διδύμου. Η χρήση της ημερήσιας κατανάλωσης κρίθηκε απαραίτητη για τη διασφάλιση μιας δίκαιης σύγκρισης, καθώς οι περίοδοι τιμολόγησης των λογαριασμών ρεύματος διέφεραν σε πλήθος ημερών.

Από πλευράς αρχιτεκτονικής ροής δεδομένων, το \textlatin{baseline} και οι πρώτες ενεργειακές χρονοσειρές συγκεντρώθηκαν στο τοπικό \textlatin{MVP} μέσω της υποδομής \textlatin{openHAB}/\textlatin{InfluxDB}. Η ίδια τηλεμετρία, αφού ωρίμασε λειτουργικά και εμπλουτίστηκε με μεγαλύτερο χρονικό ορίζοντα παρατήρησης, μεταφέρθηκε στο περιβάλλον της \textlatin{Enterprise Cloud SaaS} πλατφόρμας, όπου κατέστη δυνατή η ενοποιημένη ανάλυση, η συγκριτική εποπτεία \textlatin{baseline} και βελτιστοποιημένης λειτουργίας, καθώς και η παραγωγή προβολών μέσω \textlatin{AI Energy Forecasting} για τις μελλοντικές ταμειακές ροές και τους δείκτες \textlatin{ROI}.

\subsection{Συνολική Σύγκριση Ακατέργαστων Δεδομένων (\textlatin{Raw Comparison} -- 10,9\%)}
Η πρώτη ανάλυση αφορά το σύνολο των διαθέσιμων δεδομένων από τον Σεπτέμβριο του 2024 έως τον Νοέμβριο του 2025.
\begin{itemize}
\item \textbf{Περίοδος Πριν (Σεπ 2024 -- Φεβ 2025):} Η μέση ημερήσια κατανάλωση ανήλθε στις 201 \textlatin{kWh}/ημέρα.
\item \textbf{Περίοδος Μετά (Μάρ 2025 -- Νοέ 2025):} Η μέση ημερήσια κατανάλωση μειώθηκε στις 179 \textlatin{kWh}/ημέρα.
\end{itemize}
\textbf{Αποτέλεσμα:} Η ποσοστιαία μείωση ανέρχεται σε 10,9\%. Το νούμερο αυτό αντιπροσωπεύει τη γενική τάση εξοικονόμησης στο σύνολο του έτους, συμπεριλαμβανομένων περιόδων με χαμηλότερες απαιτήσεις θέρμανσης ή ψύξης. Η σύγκριση αυτή προκύπτει από το πλήρες σώμα των διαθέσιμων δεδομένων και αποτυπώνει την πρώτη, πιο συνολική ανάγνωση της ενεργειακής συμπεριφοράς πριν και μετά την ωρίμανση της λειτουργίας του συστήματος.
\begin{equation}
\Delta E_{\text{raw}} = \frac{201-179}{201}\cdot 100 = 10{,}9\%
\end{equation}

Αξίζει να σημειωθεί ότι ο Μάρτιος 2025, ο οποίος αντιστοιχεί στη φάση εγκατάστασης και βαθμονόμησης του συστήματος, εμφανίζει οριακή αύξηση κατανάλωσης (228 \textlatin{kWh}/ημέρα, +2,2\% \textlatin{vs baseline}). Αυτό αποδίδεται στη φυσική παρουσία τεχνικών στο πεδίο, στις δοκιμαστικές λειτουργίες εξοπλισμού και στο γεγονός ότι οι κανόνες αυτοματισμού δεν είχαν ακόμη ενεργοποιηθεί πλήρως — το σύστημα λειτουργούσε αποκλειστικά ως παθητικός παρατηρητής (\textlatin{Digital Shadow}). Η επιβεβαίωση αυτής της ερμηνείας προκύπτει άμεσα από τα δεδομένα του Απριλίου 2025, όπου μόλις ενεργοποιήθηκαν οι αυτοματισμοί, η κατανάλωση έπεσε αμέσως στις 190 \textlatin{kWh}/ημέρα (-14,8\%), αποτυπώνοντας την ομαλή καμπύλη βελτιστοποίησης που χαρακτηρίζει ένα σύστημα \textlatin{IoT} κατά τη φάση λειτουργικής ωρίμανσής του.

\subsection{Σύγκριση Περιόδου Βελτιστοποίησης (\textlatin{Weighted/Baseline Comparison} -- 19,7\%)}
Για μια βαθύτερη κατανόηση της απόδοσης του συστήματος, ορίστηκε μια Περίοδος Βάσης (\textlatin{Baseline}) που περιλαμβάνει τους μήνες αιχμής πριν την εφαρμογή του διδύμου, και μια Περίοδος Βελτιστοποίησης.
\begin{itemize}
\item \textbf{Σταθμισμένο \textlatin{Baseline}:} Λαμβάνοντας υπόψη τους μήνες με πλήρη λειτουργία του σχολείου, η κατανάλωση αναφοράς ορίστηκε στις 223 \textlatin{kWh}/ημέρα.
\item \textbf{Περίοδος Βελτιστοποίησης:} Με την πλήρη ενεργοποίηση των κανόνων του Ψηφιακού Διδύμου (όπως ο \texttt{FridgeGuardian} και η διαχείριση φωτισμού), η κατανάλωση σταθεροποιήθηκε στις 179 \textlatin{kWh}/ημέρα.
\end{itemize}
\textbf{Αποτέλεσμα:} Η μείωση αυτή της τάξεως του 19,7\% τεκμηριώνει τη λειτουργική δυναμική του συστήματος όταν αυτό λειτουργεί σε πλήρη ισχύ και διαχειρίζεται ενεργά τα φορτία του κτιρίου. Πρόκειται για το βασικό λειτουργικό \textlatin{KPI} της ενεργειακής αξιολόγησης, το οποίο δεν περιορίζεται στη συλλογή των δεδομένων από το αρχικό \textlatin{MVP}, αλλά αναδεικνύεται με σαφέστερο και επιχειρησιακά χρησιμότερο τρόπο μέσω των εργαλείων ανάλυσης της \textlatin{Cloud SaaS} πλατφόρμας.
\begin{equation}
\Delta E_{\text{weighted}} = \frac{223-179}{223}\cdot 100 = 19{,}7\%
\end{equation}

\subsection{\textlatin{Peak Performance}: Σύγκριση \textlatin{Year-over-Year} (32,5\%)}
Η πιο εντυπωσιακή απόδειξη της αξίας του Ψηφιακού Διδύμου προκύπτει από τη σύγκριση του Νοεμβρίου 2025 με τον Νοέμβριο 2024:
\begin{itemize}
\item \textbf{Νοέμβριος 2024:} 252 \textlatin{kWh}/ημέρα.
\item \textbf{Νοέμβριος 2025:} 170 \textlatin{kWh}/ημέρα.
\item \textbf{Μείωση:} 32,5\%.
\end{itemize}

Η ανάλυση των δεδομένων αναδεικνύει δύο διακριτές κορυφαίες επιδόσεις, κάθε μία με διαφορετική μεθοδολογική σημασία.

Ο Οκτώβριος 2025 αποτελεί τον μήνα με την κορυφαία απόδοση έναντι του γενικού \textlatin{Baseline}: με κατανάλωση 145 \textlatin{kWh}/ημέρα έναντι των 223 \textlatin{kWh}/ημέρα του σταθμισμένου μέσου όρου, καταγράφεται μείωση -35,0\%. Ο αριθμός αυτός αποτυπώνει την απόλυτη αποτελεσματικότητα του συστήματος κάτω από ιδανικές λειτουργικές συνθήκες.

Ο Νοέμβριος 2025 αποτελεί τον μήνα με την κορυφαία \textlatin{Year-over-Year (YoY)} απόδοση: με κατανάλωση 170 \textlatin{kWh}/ημέρα έναντι 252 \textlatin{kWh}/ημέρα τον Νοέμβριο 2024, καταγράφεται μείωση -32,5\%. Η σύγκριση αυτή έχει ιδιαίτερη ακαδημαϊκή αξία, διότι ελέγχει το σύστημα υπό παρόμοιες εποχιακές και λειτουργικές συνθήκες (ίδιος μήνας, ίδιο σχολικό πρόγραμμα, παρόμοιες καιρικές συνθήκες), απομονώνοντας έτσι την πραγματική επίδραση της πλατφόρμας από εποχιακές διακυμάνσεις. Πρόκειται για τον ισχυρότερο δείκτη αποτελεσματικότητας του Ψηφιακού Διδύμου κατά τους ενεργοβόρους χειμερινούς μήνες.
Η τάση αυτή είναι συμβατή με ευρωπαϊκά ευρήματα για ενεργειακή βελτιστοποίηση κτιρίων και δείκτες έξυπνης ετοιμότητας (\textlatin{SRI}) \cite{Apostolopoulos2022, Markoska2019, EUParliament2024, Fokaides2020, Zamanidou2024}. Τα συγκεντρωτικά ενεργειακά αποτελέσματα παρουσιάζονται στον Πίνακα~\ref{tab:energy_comparison}, ενώ η μηνιαία εξέλιξη της κατανάλωσης κατά την πιλοτική φάση αποτυπώνεται στον Πίνακα~\ref{tab:monthly_evolution}.

\begin{table}[htbp]
\centering
\caption{Συγκριτικός πίνακας ημερήσιας ενεργειακής κατανάλωσης Πριν και Μετά}
\label{tab:energy_comparison}
\begin{tabular}{|l|c|c|c|}
\hline
\textbf{Μετρική} & \textbf{Περίοδος \textlatin{Baseline}} & \textbf{Ψηφιακό Δίδυμο} & \textbf{Μεταβολή} \\
\hline
Περίοδος Μέτρησης & Σεπ. 2024 - Φεβ. 2025 & Μαρ. 2025 - Νοε. 2025 & - \\
\hline
Σύνολο Ημερών & 122 ημέρες & 226 ημέρες & - \\
\hline
Μέση Ημερήσια Κατανάλωση & \textbf{201 \textlatin{kWh}/ημέρα} & \textbf{179 \textlatin{kWh}/ημέρα} & \textbf{-10,9\%} \\
\hline
Σταθμισμένη Σύγκριση (\textlatin{Baseline}) & 223 \textlatin{kWh}/ημέρα & 179 \textlatin{kWh}/ημέρα & \textbf{-19,7\%} \\
\hline
Σύγκριση Νοεμβρίου (\textlatin{YoY}) & 252 \textlatin{kWh}/ημέρα & 170 \textlatin{kWh}/ημέρα & \textbf{-32,5\%} \\
\hline
\end{tabular}
\end{table}

Σημειώνεται ότι ο Πίνακας Β.2 του Παραρτήματος Β παρουσιάζει \textlatin{YoY} σύγκριση αποκλειστικά για τους μήνες Σεπτέμβριο--Νοέμβριο, καθώς αυτοί είναι οι μόνοι μήνες για τους οποίους υπάρχουν αντίστοιχα τιμολογημένα δεδομένα και στις δύο χρονιές (2024 και 2025). Οι μήνες Μάρτιος--Αύγουστος 2025 δεν διαθέτουν αντίστοιχο \textlatin{baseline} 2024 στα διαθέσιμα στοιχεία τιμολόγησης και επομένως εξαιρούνται από τη \textlatin{YoY} ανάλυση.

\begin{table}[htbp]
\centering
\caption{Μηνιαία εξέλιξη κατανάλωσης κατά τη διάρκεια της πιλοτικής εφαρμογής (Μάρ--Νοε 2025).}
\label{tab:monthly_evolution}
\begin{tabular}{|l|c|c|l|}
\hline
\textbf{Μήνας} & \textbf{\textlatin{kWh}/ημέρα} & \textbf{\textlatin{vs Baseline (223)}} & \textbf{Φάση Λειτουργίας} \\
\hline
Μάρτιος 2025 & 228 & +2,2\% & Εγκατάσταση \& Βαθμονόμηση \\
\hline
Απρίλιος 2025 & 190 & -14,8\% & Πρώτη Βελτιστοποίηση \\
\hline
Σεπτέμβριος 2025 & 197 & -11,7\% & Σταθερή Λειτουργία \\
\hline
Οκτώβριος 2025 & 145 & -35,0\% & Κορυφαία Απόδοση (\textlatin{vs Baseline}) \\
\hline
Νοέμβριος 2025 & 170 & -23,8\% & Κορυφαία Απόδοση (\textlatin{YoY: -32,5\%}) \\
\hline
\end{tabular}
\end{table}

\subsection*{Ερμηνεία \textlatin{YoY} Αυξήσεων (Σεπτέμβριος--Οκτώβριος)}
Στον Πίνακα Β.2 παρατηρείται ότι οι μήνες Σεπτέμβριος και Οκτώβριος 2025 εμφανίζουν \textlatin{YoY} αύξηση (+20,9\% και +8,2\% αντίστοιχα) σε σχέση με τους ίδιους μήνες του 2024. Η ερμηνεία αυτής της παρατήρησης απαιτεί διαχωρισμό των δύο μεθοδολογιών σύγκρισης.

Οι μήνες Σεπτέμβριος και Οκτώβριος 2024 παρουσίασαν ασυνήθιστα χαμηλή κατανάλωση (163 \textlatin{kWh}/ημέρα και 134 \textlatin{kWh}/ημέρα αντίστοιχα), σημαντικά κατώτερη από το σταθμισμένο \textlatin{baseline} (223 \textlatin{kWh}/ημέρα), πιθανώς λόγω ευνοϊκών καιρικών συνθηκών ή μειωμένης λειτουργίας κατά τις μεταβατικές εποχές. Συνεπώς, η \textlatin{YoY} αύξηση δεν αντανακλά μείωση αποδοτικότητας του συστήματος, αλλά απλώς υψηλότερη βάση σύγκρισης τους αντίστοιχους μήνες του 2025.

Αντίθετα, η σύγκριση έναντι του σταθμισμένου \textlatin{Baseline} επιβεβαιώνει ότι και οι δύο μήνες παρέμειναν κάτω από το επίπεδο αναφοράς (197 \textlatin{kWh}/ημέρα = -11,7\% και 145 \textlatin{kWh}/ημέρα = -35,0\% αντίστοιχα), τεκμηριώνοντας τη σταθερή αποτελεσματικότητα της πλατφόρμας. Για τον λόγο αυτό, ο κύριος δείκτης αξιολόγησης της παρούσας μελέτης (\textlatin{KPI}) παραμένει η σύγκριση \textlatin{vs Baseline} (-19,7\%), η οποία είναι ανεξάρτητη από εποχιακές διακυμάνσεις μεμονωμένων μηνών.

\section{Αποτελέσματα Τεχνοοικονομικής Ανάλυσης (\textlatin{ROI})}
Βάσει των πραγματικών δεδομένων ενεργειακής εξοικονόμησης και λαμβάνοντας υπόψη μια μέση τιμή ενέργειας της τάξης των 0,162 \euro/\textlatin{kWh}, τα οικονομικά μεγέθη της πλατφόρμας είναι εξαιρετικά βιώσιμα. Το συνολικό κόστος της αρχικής επένδυσης (\textlatin{CAPEX}), το οποίο περιλαμβάνει τον εξοπλισμό \textlatin{Edge}, τις άδειες και την ανάπτυξη του λογισμικού, υπολογίστηκε στα 19.500 \euro.

Σημειώνεται ότι το αναλυτικό \textlatin{BOM} του Παραρτήματος Α (245,17~\euro) αφορά αποκλειστικά τον πιλοτικό \textlatin{edge} εξοπλισμό της εγκατάστασης. Το συνολικό \textlatin{CAPEX} των 19.500 \euro αναφέρεται στο πλήρες έργο επίδειξης και περιλαμβάνει, πέραν του υλικού, την ανάπτυξη λογισμικού, την παραμετροποίηση της υποδομής \textlatin{openHAB}/\textlatin{MQTT}/\textlatin{InfluxDB}, την υλοποίηση της διεπαφής οπτικοποίησης, την εγκατάσταση, το \textlatin{commissioning} και την τεχνική υποστήριξη της πιλοτικής λειτουργίας.

\begin{table}[htbp]
\centering
\caption{Συμφιλίωση πιλοτικού \textlatin{BOM} και συνολικού κόστους έργου (\textlatin{CAPEX})}
\label{tab:capex_reconciliation}
\begin{tabular}{@{}p{9.2cm}r@{}}
\toprule
\textbf{Συνιστώσα Κόστους} & \textbf{Ποσό} \\
\midrule
Πιλοτικός \textlatin{edge} εξοπλισμός (\textlatin{hardware BOM}, Παράρτημα Α) & 245,17~\euro \\
Ανάπτυξη λογισμικού, παραμετροποίηση πλατφόρμας, εγκατάσταση, \textlatin{commissioning} και τεχνική υποστήριξη πιλοτικής λειτουργίας & 19.254,83~\euro \\
\midrule
\textbf{Συνολικό κόστος έργου (\textlatin{CAPEX})} & \textbf{19.500,00 \euro} \\
\bottomrule
\end{tabular}
\end{table}

Ο Πίνακας \ref{tab:capex_reconciliation} καθιστά σαφές ότι το ποσό του Παραρτήματος Α δεν αντιστοιχεί στο πλήρες κόστος του έργου, αλλά μόνο στο υλικό του πιλοτικού κόμβου. Το κύριο τμήμα του \textlatin{CAPEX} αποδίδεται στις υπηρεσίες μηχανικής υλοποίησης, ολοκλήρωσης και επιχειρησιακής θέσης σε λειτουργία της πλατφόρμας.

Στον αντίποδα, τα ετήσια καθαρά οφέλη (\textlatin{OPEX savings}), προερχόμενα από την άμεση εξοικονόμηση ηλεκτρικής ενέργειας και τη βελτιστοποίηση της προγνωστικής συντήρησης, ανέρχονται σε 5.500 \euro. Η ταμειακή αυτή ροή συνεπάγεται περίοδο απόσβεσης (\textlatin{Payback Period}) 42,5 μηνών (περίπου 3,5 έτη). Το συνολικό 5ετές καθαρό κέρδος της επένδυσης εκτιμάται στα 8.000 \euro, παρουσιάζοντας Απόδοση Επένδυσης (\textlatin{5-Year ROI}) της τάξης του 41,02\%, όπως οπτικοποιείται στο Σχήμα~\ref{fig:roi_chart}.

Από μεθοδολογική σκοπιά, το οικονομικό αποτέλεσμα βασίζεται στα πραγματικά ενεργειακά δεδομένα που συγκεντρώθηκαν από την αρχική υποδομή \textlatin{MVP}, αλλά η τελική ανάλυση και η προβολή των ταμειακών ροών πραγματοποιούνται στη νέα \textlatin{Enterprise Cloud SaaS} πλατφόρμα. Στο περιβάλλον αυτό, τα μοντέλα \textlatin{AI Energy Forecasting} αξιοποιούνται για τη δημιουργία πιο εύχρηστων προβολών κόστους, σεναρίων απόσβεσης και αναφορών \textlatin{ROI}, καθιστώντας την τεχνοοικονομική εικόνα περισσότερο κατανοητή και αξιοποιήσιμη σε διοικητικό επίπεδο. Οι αναλυτικοί δείκτες και οι επιμέρους κατηγορίες κόστους και απόδοσης παρουσιάζονται στους Πίνακες~\ref{tab:capex_breakdown}, \ref{tab:technical_kpis} και \ref{tab:roi_projection}.

\begin{table}[htbp]
\centering
\caption{Ενδεικτική ανάλυση ετήσιων λειτουργικών οφελών (\textlatin{OPEX savings})}
\label{tab:opex_breakdown}
\begin{tabular}{@{}p{9.2cm}r@{}}
\toprule
\textbf{Συνιστώσα Οφέλους} & \textbf{Ποσό / έτος} \\
\midrule
Ενεργειακή εξοικονόμηση από σταθμισμένη μείωση 44 \textlatin{kWh}/ημέρα με τιμή αναφοράς 0,162 \euro/\textlatin{kWh} ($44\times365\times0.162$) & 2.601,72 \euro \\
Εκτιμώμενη μείωση κόστους λειτουργικής/προγνωστικής συντήρησης (υπόλοιπο μοντέλου) & 2.898,28 \euro \\
\midrule
\textbf{Συνολικά ετήσια οφέλη (\textlatin{OPEX savings})} & \textbf{5.500,00 \euro} \\
\bottomrule
\end{tabular}
\end{table}

Ο Πίνακας \ref{tab:opex_breakdown} διαχωρίζει ρητά το ενεργειακό σκέλος από το σκέλος συντήρησης, ώστε η τιμή των 5.500 \euro/έτος να είναι αναπαραγώγιμη και διαφανής ως προς τις παραδοχές του μοντέλου.

\begin{figure}[htbp]
\centering
\includegraphics[width=\textwidth]{fig_6_2}
\caption{Διάγραμμα απόσβεσης της επένδυσης και προβολή ταμειακών ροών πενταετίας}
\label{fig:roi_chart}
\end{figure}

Πέραν των οικονομικών δεικτών, η αξιολόγηση της πλατφόρμας περιλαμβάνει και τεχνικούς δείκτες αξιοπιστίας. Κατά την περίοδο Μαρτίου–Νοεμβρίου 2025, το σύστημα παρουσίασε \textlatin{uptime} 99.95\%, τιμή που επιβεβαιώνει την καταλληλότητά του για επιχειρησιακή χρήση 24/7 σε εκπαιδευτικά περιβάλλοντα. Ο μέσος χρόνος διάδοσης δεδομένων από τον αισθητήρα ως την τελική διεπαφή χρήστη (\textlatin{end-to-end latency}) παρέμεινε κάτω από 5 δευτερόλεπτα, επιβεβαιώνοντας τον πραγματικού χρόνου χαρακτήρα της πλατφόρμας.

\begin{table}[htbp]
\centering
\caption{Ενδεικτικός επιμερισμός του συνολικού κόστους επένδυσης (\textlatin{CAPEX})}
\label{tab:capex_breakdown}
\begin{tabular}{|p{4.5cm}|p{6.5cm}|r|}
\hline
\textbf{Κατηγορία} & \textbf{Περιγραφή} & \textbf{Κόστος} \\
\hline
Κόστος Αισθητήρων / \textlatin{Hardware} & \textlatin{Edge \& IoT Nodes (ESP32, Shelly 1PM, Sonoff, ESP8266, DS18B20,} μαγνητικές επαφές, \textlatin{Raspberry Pi 4)} & 245,17~\euro \\
\hline
Κόστος Πλατφόρμας \& \textlatin{Cloud Integration} & \textlatin{Software stack, Supabase, openHAB} παραμετροποίηση, \textlatin{React 18 frontend} ανάπτυξη & 7.500,00 \euro \\
\hline
Κόστος Εγκατάστασης & \textlatin{Deployment}, καλωδιώσεις, δικτύωση, \textlatin{commissioning} στο πεδίο & 7.500,00 \euro \\
\hline
Κόστος Εκπαίδευσης / Παραμετροποίησης & \textlatin{Onboarding}, τεκμηρίωση, τεχνική υποστήριξη πιλοτικής λειτουργίας & 4.306,83 \euro \\
\hline
\textbf{Συνολική Επένδυση (\textlatin{CAPEX})} & & \textbf{19.500,00 \euro} \\
\hline
\end{tabular}
\end{table}

\begin{table}[htbp]
\centering
\caption{Τεχνικοί δείκτες αξιοπιστίας και απόδοσης της πλατφόρμας}
\label{tab:technical_kpis}
\begin{tabular}{|l|l|l|}
\hline
\textbf{Δείκτης} & \textbf{Τιμή} & \textbf{Μέθοδος Επαλήθευσης} \\
\hline
Ακρίβεια Θορύβου (\textlatin{INMP441}) & ±2 \textlatin{dBA} & Σύγκριση με βιομηχανικό ηχόμετρο αναφοράς \\
\hline
Ακρίβεια Θερμοκρασίας (\textlatin{DHT22}) & ±0.5°C & Σύγκριση με βαθμονομημένο θερμόμετρο \\
\hline
Ακρίβεια Υγρασίας (\textlatin{DHT22}) & ±2\% \textlatin{RH} & Σύγκριση με βαθμονομημένο υγρόμετρο \\
\hline
Συνολική Αξιοπιστία Δεδομένων & 98.5\% & Διασταύρωση με όργανα αναφοράς \\
\hline
\textlatin{System Uptime} & 99.95\% & Συνεχής καταγραφή Μαρ–Νοε 2025 \\
\hline
\textlatin{End-to-End Latency} & < 5 \textlatin{sec} & Μέτρηση αισθητήρας $\rightarrow$ \textlatin{UI} \\
\hline
\textlatin{Sensor} $\rightarrow$ \textlatin{MQTT} & $\sim$45 \textlatin{ms} & \textlatin{Network timing analysis} \\
\hline
\textlatin{MQTT} $\rightarrow$ \textlatin{openHAB} & $\sim$28 \textlatin{ms} & \textlatin{openHAB log analysis} \\
\hline
\textlatin{openHAB} $\rightarrow$ \textlatin{UI Update} & $\sim$19 \textlatin{ms} & \textlatin{Browser DevTools} \\
\hline
\end{tabular}
\end{table}

\noindent\textit{Σημείωση: Η κατανομή ανά κατηγορία είναι ενδεικτική. Το συνολικό \textlatin{CAPEX} των 19.500~\euro\ αποτελεί το μέγεθος αναφοράς της τεχνοοικονομικής ανάλυσης.}

\begin{table}[htbp]
\centering
\caption{Εξέλιξη ετήσιων λειτουργικών οφελών και συνοπτικοί δείκτες απόδοσης επένδυσης}
\label{tab:roi_projection}
\begin{tabular}{|l|l|r|}
\hline
\textbf{Περίοδος / Δείκτης} & \textbf{Περιγραφή} & \textbf{Τιμή} \\
\hline
\textlatin{Year 1} & Ετήσια λειτουργικά οφέλη (\textlatin{OPEX Savings}) & 5.500,00 \euro \\
\hline
\textlatin{Year 2} & Σωρευτικά λειτουργικά οφέλη & 11.000,00 \euro \\
\hline
\textlatin{Year 3} & Σωρευτικά λειτουργικά οφέλη & 16.500,00 \euro \\
\hline
\textlatin{Year 5} & Σωρευτικά λειτουργικά οφέλη & 27.500,00 \euro \\
\hline
5-ετές Καθαρό Κέρδος & Μετά την αφαίρεση της αρχικής επένδυσης & 8.000,00 \euro \\
\hline
Περίοδος Απόσβεσης & \textlatin{Payback Period} & 42,5 μήνες \\
\hline
Απόδοση Επένδυσης 5ετίας & \textlatin{5-year ROI} & 41,02\% \\
\hline
\end{tabular}
\end{table}




